Let $\mathcal X=\{x_i\}_{i=0}^{N-1}\subset\mathbb R^d$ be a database
with stable identifiers. For query $q$, the exact task selects the $k$
smallest raw-vector distances
\begin{equation}
 \rawdist(q,x)=\|q-x\|_2^2.
 \label{eq:2-raw}
\end{equation}
Ties require a specified ordering or a set of equally valid answers.
This thesis evaluates the stored GT top-10 set, with its fixed tie choices.
ANN reduces the cost of accessing and scoring vectors while permitting
missed members of that set. The primary quality measure is Recall@10
(\cref{eq:3-recall}); it is an empirical set-recovery measure. A high measured
recall does not establish a worst-case multiplicative distance guarantee.

There are four relevant cost axes. Query throughput counts queries completed
per second at a stated batch and timing boundary. Batch latency counts
elapsed time for that call; throughput alone supplies no concurrent-service
tail guarantee. Construction cost includes the specified training and
encoding stages. Storage distinguishes persistent bytes per vector from
model, workspace, context and temporary allocation. Chapter 3 defines the
actual measurement boundaries. A method may improve one axis while worsening
another; equal recall, equal memory and equal construction effort are
different comparisons.

Tuning produces a finite set of operating points. Its Pareto frontier keeps
points whose recall and throughput are not simultaneously improved by
another measured configuration. At a target recall, the fastest qualifying
point is an available configuration, while interpolation estimates between
neighbouring points. An uncovered region supplies no evidence of either
dominance or impossibility. These distinctions are especially relevant when
the compared systems use different code precisions and graph/index sizes.
The detailed frontier estimators and their endpoint tolerance are fixed in
\cref{sec:3-4}.
